\documentclass[11pt,a4paper]{article}

\usepackage[utf8]{inputenc}
\usepackage[T1]{fontenc}
\usepackage{amsmath,amssymb}
\usepackage{booktabs}
\usepackage{array,tabularx}
\usepackage{graphicx}
\usepackage{hyperref}
\usepackage[margin=2.5cm]{geometry}
\usepackage{algorithm}
\usepackage{algorithmic}
\usepackage{xcolor}
\usepackage{enumitem}

\definecolor{accent}{RGB}{37,99,235}
\hypersetup{
    hypertexnames=false,
    colorlinks=true,
    linkcolor=accent,
    citecolor=accent,
    urlcolor=accent
}

\title{\textbf{Two-Stage Hierarchical Postprocessing\\for SciScape Paper Clustering}}
\author{Young Jin Kim}
\date{April 2026}

\begin{document}
\maketitle

\begin{abstract}
SciScape builds a publication-level research landscape by repeatedly running
CPM Leiden clustering, resolving undersized clusters, and contracting the
result into the next hierarchy level. This report documents an internal,
opt-in extension that adds a second postprocessing stage at every hierarchy
level: after the existing small-cluster repair, the system attempts to reduce
oversized clusters using split--repair probes and boundary-node trimming. The
method preserves the default product behavior, prioritizes CPM quality by
default, and records explicit diagnostics whenever a hard maximum-size cap is
requested but cannot be satisfied. The design is motivated by a recurring
hierarchical failure mode: small-cluster repair can improve minimum-size
constraints while leaving or creating large clusters that dominate subsequent
graph contraction. The proposed automation makes the level transition
``raw Leiden $\rightarrow$ small repair $\rightarrow$ oversize repair
$\rightarrow$ next hierarchy'' explicit, cache-aware, and auditable.
\end{abstract}

%% ----------------------------------------------------------------------
\section{Problem Context}

SciScape represents papers as nodes in a weighted undirected graph. Edges come
from citation and semantic similarity layers, then CPM Leiden is used to find
communities. The hierarchy builder runs this process level by level:

\begin{enumerate}[nosep]
    \item run Leiden at the current level;
    \item repair clusters below the level's minimum document weight;
    \item project membership back to original papers;
    \item contract clusters into super-nodes for the next level.
\end{enumerate}

This is a natural architecture for multi-scale scientific mapping, but it has
an asymmetry: the existing postprocess stage enforces a lower size constraint,
while the upper size target is used mainly during resolution selection. Once a
small-cluster repair merges low-weight groups into larger neighbors, the
largest cluster can remain above target. If this cluster is contracted as a
single super-node, it becomes a coarse unit in the next level and can distort
the hierarchy.

\paragraph{Core motivation.}
The target maximum cluster size is not just a reporting statistic. In a
hierarchy, it controls the granularity of the next graph. A level that contains
one very large cluster and many moderate clusters passes an imbalanced
state into contraction. This can make upper levels less interpretable even
when the current level has high CPM quality.

%% ----------------------------------------------------------------------
\section{Theoretical Grounding}

\subsection{CPM Objective}

SciScape uses the Constant Potts Model (CPM) quality function. For a partition
$\mathcal{P}$ and resolution $\gamma$, CPM can be written as:

\begin{equation}
    Q_{\text{CPM}}(\mathcal{P}; \gamma)
    =
    \sum_{C \in \mathcal{P}}
    \left[
        e_C - \gamma \binom{|C|}{2}
    \right],
\end{equation}

where $e_C$ is the total internal edge weight in cluster $C$. CPM is used
because it is designed to avoid the classical modularity resolution limit
described by Fortunato and Barthelemy~\cite{fortunato2007resolution}, and
formalized in the resolution-limit-free discussion by Traag, Van Dooren, and
Nesterov~\cite{traag2011narrow}. Leiden is used as the optimizer because it
improves on Louvain by guaranteeing well-connected communities under its
refinement procedure~\cite{traag2019leiden}.

\subsection{Hierarchy-Specific Size Semantics}

At hierarchy level $\ell$, the graph nodes may represent either original papers
or contracted paper groups. Let $s_i$ denote the document weight of current
node $i$. For a cluster $C$, its document weight is:

\begin{equation}
    W(C) = \sum_{i \in C} s_i.
\end{equation}

The lower and upper targets at level $\ell$ are:

\begin{align}
    W_{\min}^{(\ell)} &= \texttt{min\_sizes[level]}, \\
    W_{\max}^{(\ell)} &= \frac{\texttt{targets[level]}}{100}
                         \sum_{i \in V_\ell} s_i.
\end{align}

For the first level, $s_i=1$ for every paper. For contracted levels, $s_i$ is
the number of original papers represented by super-node $i$.

%% ----------------------------------------------------------------------
\section{Method Overview}

The new method is a two-stage postprocess applied at each hierarchy level.
It is disabled by default and is activated only through an internal
\texttt{HierarchyPostprocessConfig}.

\begin{algorithm}[h]
\caption{Level-wise two-stage hierarchy postprocess}
\begin{algorithmic}[1]
\REQUIRE Current graph $G_\ell$, Leiden membership $m_{\text{raw}}$, level targets
\STATE $m_{\text{small}} \leftarrow \textsc{SmallClusterRepair}(m_{\text{raw}}, W_{\min})$
\IF{hierarchy postprocess is disabled}
    \STATE \RETURN $m_{\text{small}}$
\ENDIF
\STATE $m \leftarrow m_{\text{small}}$
\FOR{$t = 1$ to \texttt{apply\_iterations}}
    \STATE $O \leftarrow \{C : W(C) > W_{\max}\}$
    \IF{$O = \emptyset$}
        \STATE \textbf{break}
    \ENDIF
    \STATE $P \leftarrow \textsc{SplitRepairProbes}(G_\ell, m, O, \gamma)$
    \STATE $S \leftarrow \textsc{SelectNonConflictingCandidates}(P)$
    \STATE $m' \leftarrow \textsc{ApplySplitRepair}(G_\ell, m, S)$
    \IF{$Q_{\text{CPM}}(m') - Q_{\text{CPM}}(m) < 0$}
        \STATE \textbf{break}
    \ENDIF
    \STATE $m \leftarrow m'$
\ENDFOR
\STATE $m_{\text{trim}} \leftarrow \textsc{BoundaryTrim}(G_\ell, m, W_{\max})$
\STATE accept or reject $m_{\text{trim}}$ according to policy
\RETURN accepted membership, or $m_{\text{small}}$ as fallback
\end{algorithmic}
\end{algorithm}

\subsection{Stage 1: Small-Cluster Repair}

The first stage is the existing hierarchy behavior. It resolves clusters below
$W_{\min}$ using the Rust postprocess backend. On contracted graphs, the repair
uses node weights, so the threshold is interpreted as original-paper count
rather than current super-node count.

This stage remains mandatory because tiny clusters and singletons often arise
from CPM resolution search and graph sparsity. Without this repair, the finest
level can become fragmented and upper levels inherit many low-weight
super-nodes.

\subsection{Stage 2A: Oversize Split--Repair}

After small-cluster repair, the method identifies current oversize clusters:

\begin{equation}
    O = \{C \in \mathcal{P}: W(C) > W_{\max}\}.
\end{equation}

For each candidate in $O$, split--repair probes force a higher-resolution split
inside the source cluster and then repair the result back at the baseline
resolution. Each probe records:

\begin{itemize}[nosep]
    \item source cluster document weight;
    \item induced subgraph size;
    \item number of proposed parts;
    \item escaped source weight;
    \item retained source units;
    \item predicted and exact CPM quality deltas;
    \item residual oversize after the candidate move.
\end{itemize}

Selection defaults to \texttt{oversize\_first}. This ranking is intentionally
different from a pure utility-per-cost ranking: for hierarchy construction,
the immediate objective is to reduce clusters that exceed the level-local
maximum while preserving CPM quality.

\subsection{Stage 2B: Oversize Boundary Trim}

Split--repair can fail when a cluster has no stable multi-core split, or when
all split proposals restore the original source cluster after repair. The trim
stage therefore tries a smaller operation: moving boundary nodes out of
oversize clusters into neighboring clusters, subject to receiver capacity and
a per-move $\Delta Q$ lower bound.

Let $\delta q_v$ be the predicted CPM quality change of moving boundary node
$v$. The trim stage accepts moves satisfying:

\begin{equation}
    \delta q_v \geq \tau_{\text{trim}},
\end{equation}

and additionally prevents the target cluster from exceeding $W_{\max}$. A
quality-floor prefix is used so that the committed prefix of trim moves keeps
the final exact CPM quality above the policy floor.

%% ----------------------------------------------------------------------
\section{Acceptance Policies}

The extension supports two policies. The default is \texttt{quality\_first}.

\begin{table}[h]
\centering
\caption{Oversize postprocess acceptance policies}
\begin{tabular}{lll}
\toprule
Policy & Trim default & Final acceptance rule \\
\midrule
\texttt{quality\_first}
    & $\tau_{\text{trim}} = 0$
    & $\Delta Q_{\text{exact}} \geq \Delta Q_{\min}$ \\
\texttt{hard\_cap}
    & $\tau_{\text{trim}} = -1$
    & $\Delta Q_{\text{exact}} \geq \Delta Q_{\min}$ and
      $\max_C W(C) \leq W_{\max}$ \\
\bottomrule
\end{tabular}
\end{table}

\paragraph{\texttt{quality\_first}.}
This is the operating default. It accepts a postprocessed membership if final
exact CPM quality is not below the configured floor. It does not require the
maximum-size target to be fully satisfied. The summary records whether the cap
was satisfied as a diagnostic field.

\paragraph{\texttt{hard\_cap}.}
This is an internal experimental policy for cases where the maximum document
weight must be enforced. It allows slightly negative trim moves by default, but
only accepts the result if the final exact quality floor and the hard maximum
condition are both satisfied. If the hard cap is not satisfied, the diagnostic
membership can be written, but the next hierarchy level uses the safer
$m_{\text{small}}$ fallback.

%% ----------------------------------------------------------------------
\section{Cache and Artifacts}

Default behavior is unchanged when hierarchy postprocess is disabled. Existing
level caches remain backward compatible. When enabled, each level writes
additional artifacts:

\begin{itemize}[nosep]
    \item \texttt{cache\_dir/<level>/postprocess/summary.json}
    \item \texttt{cache\_dir/<level>/postprocess/oversize\_boundary\_trim\_moves.csv}
    \item optional \texttt{diagnostic\_membership.parquet} for rejected hard-cap
          or quality-floor cases with proposed changes.
\end{itemize}

The level \texttt{meta.json} is extended only for enabled runs:

\begin{itemize}[nosep]
    \item \texttt{hierarchy\_postprocess\_enabled}
    \item \texttt{postprocess\_status}
    \item \texttt{oversize\_policy}
    \item \texttt{target\_min\_doc\_weight}
    \item \texttt{target\_max\_doc\_weight}
    \item \texttt{target\_max\_satisfied}
    \item \texttt{postprocess\_quality\_delta}
    \item \texttt{postprocess\_config\_hash}
\end{itemize}

When enabled, the hierarchy cache is valid only if the stored
\texttt{postprocess\_config\_hash} matches the active config hash. This avoids
silently reusing memberships generated under a different postprocess policy.

%% ----------------------------------------------------------------------
\section{Data and Evidence}

\subsection{Design Evidence from Field-Level Experiments}

The broader SciScape pipeline has been evaluated on OpenAlex fields with
contrasting citation densities. The multi-layer report documents two important
observations that motivate hierarchy-level automation:

\begin{table}[h]
\centering
\caption{Representative field characteristics from previous SciScape experiments}
\begin{tabularx}{\textwidth}{p{0.30\textwidth}p{0.16\textwidth}X}
\toprule
Dataset & Scale & Key observation \\
\midrule
Field 15: Chemical Engineering
    & 115K papers
    & Citation layers have high coverage; auto-$\gamma$ selects balanced clusters. \\
Field 12: Arts and Humanities
    & 754K papers
    & Sparse citation coverage makes embedding edges important for connectivity. \\
\bottomrule
\end{tabularx}
\end{table}

These datasets show that the same hierarchy code must handle both dense
scientific fields and sparse humanities fields. A fixed postprocess assumption
is therefore brittle: after contraction, cluster document weights can vary
substantially even when resolution search is successful at the current level.

\subsection{Implementation Verification}

The current implementation was validated with unit, integration, and Rust
backend tests.

\begin{table}[h]
\centering
\caption{Verification commands run after implementation}
\begin{tabularx}{\textwidth}{p{0.30\textwidth}X}
\toprule
Suite & Result \\
\midrule
Rust backend
    & \texttt{cargo test --manifest-path rust/Cargo.toml}: unit and integration tests passed. \\
Targeted Python
    & Hierarchy, adaptive-refinement, and Rust-wrapper tests: 77 passed, 1 warning. \\
Full Python
    & \texttt{uv run python -m pytest -q}: 987 passed, 1218 warnings. \\
\bottomrule
\end{tabularx}
\end{table}

The added tests cover:

\begin{itemize}[nosep]
    \item config defaults: disabled and \texttt{quality\_first};
    \item target maximum calculation from level percentage and current document weight;
    \item policy-specific trim bounds;
    \item enabled artifact and metadata creation;
    \item cache invalidation on config-hash mismatch;
    \item hard-cap rejection and fallback to the small-postprocessed membership.
\end{itemize}

%% ----------------------------------------------------------------------
\section{Operational Interpretation}

\subsection{Default Product Behavior}

The method is intentionally not exposed as a public CLI flag in this stage.
It is an internal development feature. Without an explicit config, the product
pipeline behaves exactly as before:

\begin{equation}
    m_{\text{level}} = m_{\text{small}}.
\end{equation}

\subsection{Recommended Internal Use}

For exploratory hierarchy-building, use:

\begin{verbatim}
HierarchyPostprocessConfig(
    enabled=True,
    oversize_policy="quality_first"
)
\end{verbatim}

Use \texttt{hard\_cap} only for experiments where the maximum document-weight
constraint is the primary objective and quality loss must be explicitly
measured.

\subsection{Interpretability of Outcomes}

The postprocess status should be interpreted as follows:

\begin{table}[h]
\centering
\caption{Status semantics}
\begin{tabularx}{\textwidth}{p{0.42\textwidth}X}
\toprule
Status & Meaning \\
\midrule
{\small\texttt{no\_current\_oversize\_candidates}}
    & Small-repaired membership already has no cluster above target. \\
\texttt{committed}
    & Oversize postprocess was accepted and used for projection/contraction. \\
\texttt{quality\_below\_floor}
    & Proposed membership was rejected because exact CPM quality fell below floor. \\
\texttt{hard\_cap\_not\_satisfied}
    & Diagnostic was written, but fallback membership is used for next level. \\
\bottomrule
\end{tabularx}
\end{table}

%% ----------------------------------------------------------------------
\section{Limitations and Open Questions}

\begin{enumerate}[nosep]
    \item \textbf{No new Rust kernel in this stage.} The method composes
          existing Rust primitives rather than adding a single specialized
          hierarchy postprocess kernel.
    \item \textbf{Quality-first may leave oversized clusters.} This is a
          deliberate default. The cap satisfaction flag must be inspected in
          summaries.
    \item \textbf{Hard-cap feasibility is graph-dependent.} If no acceptable
          boundary moves exist, the system records diagnostics and falls back.
    \item \textbf{Runtime cost grows with oversize candidates.} The method
          runs split--repair probes repeatedly, so it should remain opt-in
          until larger field-level runtime profiles are collected.
    \item \textbf{Semantic quality still needs external validation.} Structural
          balance is not the same as topic coherence. LLM or expert review
          should be used for representative boundary cases.
\end{enumerate}

%% ----------------------------------------------------------------------
\section{Conclusion}

The two-stage hierarchy postprocess makes a previously implicit need explicit:
after repairing clusters below the minimum size, hierarchy construction also
needs a controlled mechanism for diagnosing and reducing clusters above the
level target. The design preserves the CPM quality objective by default,
supports a stricter hard-cap mode for experiments, and writes enough metadata
to make every accepted or rejected level transition auditable. Most
importantly, accepted memberships are used consistently for projection,
per-level artifacts, metadata, and next-level graph contraction, so the
hierarchy reflects the actual postprocessed partition rather than an
intermediate diagnostic.

\begin{thebibliography}{9}

\bibitem{fortunato2007resolution}
S. Fortunato and M. Barthelemy.
\newblock Resolution limit in community detection.
\newblock \emph{Proceedings of the National Academy of Sciences}, 104(1):36--41, 2007.
\newblock \url{https://doi.org/10.1073/pnas.0605965104}.

\bibitem{traag2011narrow}
V. A. Traag, P. Van Dooren, and Y. Nesterov.
\newblock Narrow scope for resolution-limit-free community detection.
\newblock \emph{Physical Review E}, 84:016114, 2011.
\newblock \url{https://doi.org/10.1103/PhysRevE.84.016114}.

\bibitem{traag2019leiden}
V. A. Traag, L. Waltman, and N. J. van Eck.
\newblock From Louvain to Leiden: guaranteeing well-connected communities.
\newblock \emph{Scientific Reports}, 9:5233, 2019.
\newblock \url{https://doi.org/10.1038/s41598-019-41695-z}.

\end{thebibliography}

\end{document}
